\subsection{Schwartz 类与缓增分布}
\label{sec:ch1-schwartz-tempered-distributions}

若函数 $f$ 无限次可微, 且它的所有导数在无穷远处都迅速衰减, 则称 $f$ 属于 Schwartz 类 $\mathcal S(\mathbb R^n)$. 也就是说, 对所有 $\alpha,\beta\in\mathbb N^n$,
\[
 \sup_x|x^\alpha D^\beta f(x)|=p_{\alpha,\beta}(f)<\infty.
\]
$C_c^\infty$ 中的函数属于 $\mathcal S$, 像 $e^{-|x|^2}$ 这样不具有紧支集的函数也属于 $\mathcal S$. 集合 $\{p_{\alpha,\beta}\}$ 是 $\mathcal S$ 上的可数半范数族, 可用它在 $\mathcal S$ 上定义一个拓扑: 序列 $\{\phi_k\}$ 收敛到 $0$, 当且仅当对所有 $\alpha,\beta\in\mathbb N^n$,
\[
 \lim_{k\to\infty}p_{\alpha,\beta}(\phi_k)=0.
\]

\hypertarget{chapter-01-page-025}{}
在此拓扑下, $\mathcal S$ 是 Fréchet 空间(完备且可度量), 并且当 $1\le p<\infty$ 时在 $L^p(\mathbb R^n)$ 中稠密. 特别地, $\mathcal S\subset L^1$, 因而 \eqref{eq:ch1-fourier-transform-definition} 定义了 $\mathcal S$ 中函数的 Fourier 变换.

$\mathcal S$ 上的有界线性泛函空间 $\mathcal S'$ 称为缓增分布空间. 从 $\mathcal S$ 到 $\mathbb C$ 的线性映射 $T$ 属于 $\mathcal S'$, 当且仅当
\[
 \lim_{k\to\infty}T(\phi_k)=0
 \quad\text{只要}\quad
 \lim_{k\to\infty}\phi_k=0\text{ 于 }\mathcal S.
\]

\begin{Theorem}
\label{thm:ch1-schwartz-fourier-inversion}
Fourier 变换是从 $\mathcal S$ 到 $\mathcal S$ 的连续映射, 并满足
\begin{equation}
\label{eq:ch1-fourier-pairing-symmetry}
 \int_{\mathbb R^n}f\widehat g=\int_{\mathbb R^n}\widehat f g,
\end{equation}
以及
\begin{equation}
\label{eq:ch1-fourier-inversion}
 f(x)=\int_{\mathbb R^n}\widehat f(\xi)e^{2\pi ix\cdot\xi}\,d\xi.
\end{equation}
\end{Theorem}

等式 \eqref{eq:ch1-fourier-inversion} 称为反演公式.

为证明定理~\ref{thm:ch1-schwartz-fourier-inversion}, 需要计算一个特殊函数的 Fourier 变换.

\begin{Lemma}
\label{lem:ch1-gaussian-fourier-transform}
若 $f(x)=e^{-\pi|x|^2}$, 则 $\widehat f(\xi)=e^{-\pi|\xi|^2}$.
\end{Lemma}

\begin{Proof}
可以在 $\mathbb C$ 中积分来直接证明这一结果, 但这里给出另一种证明. 只需证明一维情形, 因为在 $\mathbb R^n$ 中, $\widehat f$ 是 $n$ 个相同积分的乘积.

函数 $f(x)=e^{-\pi x^2}$ 是微分方程
\[
 u'+2\pi xu=0,
 \qquad u(0)=1
\]
的解. 由 \eqref{eq:ch1-fourier-derivative} 和 \eqref{eq:ch1-fourier-coordinate-multiplication}, $\widehat u$ 满足同一微分方程, 且初值为
\[
 \widehat u(0)=\int_{\mathbb R}u(x)\,dx
 =\int_{\mathbb R}e^{-\pi x^2}\,dx=1.
\]
因此由唯一性, $\widehat f=f$.
\end{Proof}

\begin{Proof}
这里证明定理~\ref{thm:ch1-schwartz-fourier-inversion}. 由 \eqref{eq:ch1-fourier-derivative} 和 \eqref{eq:ch1-fourier-coordinate-multiplication},
\[
 \xi^\alpha D^\beta\widehat f(\xi)
 =C\widehat{(D^\alpha x^\beta f)}(\xi),
\]
所以
\[
 |\xi^\alpha D^\beta\widehat f(\xi)|
 \le C\|D^\alpha x^\beta f\|_1.
\]
右端的 $L^1$ 范数可由 $f$ 的有限个半范数的线性组合控制, 这说明 Fourier 变换是从 $\mathcal S$ 到自身的连续映射.

\hypertarget{chapter-01-page-026}{}
等式 \eqref{eq:ch1-fourier-pairing-symmetry} 是 Fubini 定理的直接推论, 因为 $f(x)g(y)$ 在 $\mathbb R^n\times\mathbb R^n$ 上可积.

由 \eqref{eq:ch1-fourier-dilation} 和 \eqref{eq:ch1-fourier-pairing-symmetry}, 得
\[
 \int f(x)\widehat g(\lambda x)\,dx
 =\int\widehat f(x)\lambda^{-n}g(\lambda^{-1}x)\,dx.
\]
在第一个积分中作变量代换 $\lambda x=y$, 得
\[
 \int f(\lambda^{-1}x)\widehat g(x)\,dx
 =\int\widehat f(x)g(\lambda^{-1}x)\,dx;
\]
再令 $\lambda\to\infty$, 得
\[
 f(0)\int\widehat g(x)\,dx=g(0)\int\widehat f(x)\,dx.
\]
令 $g(x)=e^{-\pi|x|^2}$. 由引理~\ref{lem:ch1-gaussian-fourier-transform},
\[
 f(0)=\int\widehat f(\xi)\,d\xi,
\]
这就是 $x=0$ 时的 \eqref{eq:ch1-fourier-inversion}. 将 $f$ 换成 $\tau_xf$, 再由 \eqref{eq:ch1-fourier-translation-modulation},
\[
 f(x)=(\tau_xf)(0)=\int\widehat{(\tau_xf)}(\xi)\,d\xi
 =\int\widehat f(\xi)e^{2\pi ix\cdot\xi}\,d\xi.
\]
\end{Proof}

令 $\widetilde f(x)=f(-x)$, 得到下述推论.

\begin{Corollary}
\label{cor:ch1-fourier-period-four}
对 $f\in\mathcal S$, 有 $\widehat{\widehat f}=\widetilde f$; 因而 Fourier 变换的周期为 $4$(即连续应用四次得到恒等算子).
\end{Corollary}

\begin{Definition}
\label{def:ch1-tempered-distribution-fourier-transform}
$T\in\mathcal S'$ 的 Fourier 变换是由
\[
 \widehat T(f)=T(\widehat f),\qquad f\in\mathcal S
\]
给出的缓增分布 $\widehat T$.
\end{Definition}

由定理~\ref{thm:ch1-schwartz-fourier-inversion}, $\widehat T$ 是缓增分布. 特别地, 若 $T$ 是可积函数, 则 $\widehat T$ 与等式 \eqref{eq:ch1-fourier-transform-definition} 定义的 Fourier 变换一致. 同样, 若 $\mu$ 是有限 Borel 测度(即 $C_0(\mathbb R^n)$ 上的有界线性泛函, 这里 $C_0(\mathbb R^n)$ 是在无穷远处消失的连续函数空间), 则 $\widehat\mu$ 是由
\[
 \widehat\mu(\xi)=\int_{\mathbb R^n}e^{-2\pi ix\cdot\xi}\,d\mu(x)
\]
给出的有界连续函数. 对原点处的 Dirac 测度 $\delta$, 这给出 $\widehat\delta=1$.

\hypertarget{chapter-01-page-027}{}
\begin{Theorem}
\label{thm:ch1-fourier-isomorphism-sprime}
Fourier 变换是从 $\mathcal S'$ 到 $\mathcal S'$ 的有界线性双射, 其逆映射也有界.
\end{Theorem}

\begin{Proof}
若 $T_n\to T$ 于 $\mathcal S'$, 则对任意 $f\in\mathcal S$,
\[
 \widehat T_n(f)=T_n(\widehat f)\longrightarrow T(\widehat f)=\widehat T(f).
\]
此外, Fourier 变换的周期为 $4$, 所以其逆映射等价于连续应用 Fourier 变换三次; 因此其逆映射也连续.
\end{Proof}

若定义 $\widetilde T$ 为 $\widetilde T(f)=T(\widetilde f)$, 则由推论~\ref{cor:ch1-fourier-period-four} 有 $\widetilde{(\widehat{\widetilde T})}=T$. 若 $\widehat T\in L^1$, 则由反演公式,
\[
 T(x)=\int_{\mathbb R^n}\widehat T(\xi)e^{2\pi ix\cdot\xi}\,d\xi;
\]
特别地, $T$ 是有界连续函数.

